\documentclass[a4paper,10pt]{article}
\usepackage[utf8]{inputenc}
\usepackage[T1]{fontenc}
\usepackage{amsmath, amssymb}
\usepackage{geometry}
\usepackage{multicol}
\geometry{margin=1.5cm}
\setlength{\parindent}{0pt}

\title{\textbf{Fiche de Révision -- Traitement du Signal}}
\date{}
\author{Raphaël Jamann}

\begin{document}
    \maketitle

    \begin{multicols}{2}[
        \section{Signaux de base}
    ]

        \textbf{Signal :} fonction représentant une grandeur physique dépendant du temps.

        \textbf{Types :}
        \begin{itemize}
            \item Continu $x(t)$
            \item Discret $x[n]$
        \end{itemize}

        \textbf{Signaux élémentaires :}
        \begin{itemize}
            \item Impulsion de Dirac $\delta(t)$ :
 
            \begin{itemize}
                \item $\int_{-\infty}^\infty \delta(t)dt = 1$ \\
                \item fonction nulle partout sauf en $t=0$
                \item $\delta(t) = \frac{d}{dt}u(t)$
                \item Propriété d'échantillonnage :
                \[
                \int x(t)\delta(t-t_0)dt = x(t_0)
                \]
            \end{itemize}

            \item Échelon unité $u(t)$ :
            \[
                u(t)=
                \begin{cases}
                1 & t \ge 0 \\
                0 & t < 0
                \end{cases}
            \]
            \item fonction rectangle : 
            \[
                \text{rect}(t) =
                \begin{cases}
                1 & |t| \le \frac{1}{2} \\
                0 & \text{sinon}
                \end{cases}
            \]
            \item Sinusoïde : $A \cos(\omega t + \phi)$
        \end{itemize}

        \textbf{Propriétés :}
        \begin{itemize}
            \item Énergie : $E = \int |x(t)|^2 dt$
            \item Puissance : $P = \lim_{T \to \infty} \frac{1}{T} \int |x(t)|^2 dt$
        \end{itemize}

        \subsection*{Opérations élémentaires sur les signaux}

        \begin{itemize}
            \item \textbf{Décalage temporel :} $x(t-t_0)$ (retard), $x(t+t_0)$ (avance)
            \item \textbf{Changement d'échelle :} $x(at)$ 
            \begin{itemize}
                \item $|a|>1$ : compression
                \item $|a|<1$ : dilatation
            \end{itemize}
            \item \textbf{Retournement :} $x(-t)$
        \end{itemize}

        \subsection*{Représentation complexe des sinusoïdes}

        \[
        e^{j2\pi f_0 t} = \cos(2\pi f_0 t) + j \sin(2\pi f_0 t)
        \]

        \[
        \cos(2\pi f_0 t) = \frac{1}{2}(e^{j2\pi f_0 t} + e^{-j2\pi f_0 t})
        \]

        \[
        \sin(2\pi f_0 t) = \frac{1}{2j}(e^{j2\pi f_0 t} - e^{-j2\pi f_0 t})
        \]


        \subsection*{Série de Fourier (signal périodique)}

        \[
        x(t) = a_0 + \sum_{k=1}^{+\infty} A_k \cos(2\pi k f_0 t + \Phi_k)
        \]

        \textbf{Coefficients :}
        \[
        a_k = \frac{2}{T_0} \int_0^{T_0} x(t)\cos(2\pi k f_0 t)\,dt
        \quad
        b_k = \frac{2}{T_0} \int_0^{T_0} x(t)\sin(2\pi k f_0 t)\,dt
        \]

        \[
        A_k = \sqrt{a_k^2 + b_k^2}, \quad \Phi_k = \arctan\left(\frac{b_k}{a_k}\right)
        \]

        \textbf{Forme complexe :}
        \[
        x(t) = \sum_{k=-\infty}^{+\infty} c_k e^{j2\pi k f_0 t}
        \]

        \[
        c_k = \frac{1}{T_0} \int_0^{T_0} x(t)e^{-j2\pi k f_0 t} dt
        \]

        \subsection*{Notions fréquentielles}

        \begin{itemize}
            \item $f_0$ : fréquence fondamentale
            \item $k f_0$ : harmoniques
            \item Spectre d'amplitude : $A_k$
            \item Spectre de phase : $\Phi_k$
        \end{itemize}


        
        \subsection*{Lien temps / fréquence}

        \begin{itemize}
            \item Signal périodique $\Rightarrow$ spectre discret
            \item Signal non périodique $\Rightarrow$ spectre continu
        \end{itemize}
            
    \end{multicols}

    \clearpage

    \begin{multicols}{2}[ 
        \section{Systèmes linéaires invariants (SLI)}
        ]

        \textbf{Définition :}
        \begin{itemize}
            \item Linéarité : superposition
            \item Invariance temporelle : décalage entrée $\Rightarrow$ décalage sortie
        \end{itemize}

        \textbf{Réponse impulsionnelle :} $h(t)$

        \textbf{Sortie d'un SLI : convolution}
        \[
        y(t) = x(t) * h(t) = \int x(\tau) h(t-\tau) d\tau
        \]

        \textbf{Propriétés de la convolution :}
        \begin{itemize}
            \item Commutativité
            \item Associativité
            \item Distributivité
        \end{itemize}

    \end{multicols}

    \begin{multicols}{2}[
    \section{Transformée de Fourier}
    ]
        \textbf{Transformée de Fourier :}
        \[
        X(f) = \int_{-\infty}^{+\infty} x(t) e^{-j2\pi ft} dt
        \]

        \textbf{Transformée inverse :}
        \[
        x(t) = \int_{-\infty}^{+\infty} X(f) e^{j2\pi ft} df
        \]

        \textbf{Interprétation :} passage du domaine temporel au domaine fréquentiel.

        \textbf{Propriétés importantes :}

        \begin{itemize}
            \item \textbf{Linéarité :}
            \[
            a x_1(t) + b x_2(t) \longleftrightarrow a X_1(f) + b X_2(f)
            \]
            \textit{Signification :} la transformée conserve les combinaisons linéaires.  
            \textit{Exemple :} somme de deux sinusoïdes $\Rightarrow$ somme de deux pics dans le spectre.

            \item \textbf{Translation temporelle :}
            \[
            x(t - t_0) \longleftrightarrow X(f)e^{-j2\pi f t_0}
            \]
            \textit{Signification :} un retard dans le temps ajoute une phase en fréquence.  
            \textit{Exemple :} décaler un signal ne change pas son amplitude fréquentielle mais modifie sa phase.

            \item \textbf{Modulation (translation fréquentielle) :}
            \[
            x(t)e^{j2\pi f_0 t} \longleftrightarrow X(f - f_0)
            \]
            \textit{Signification :} multiplication par une exponentielle $\Rightarrow$ décalage du spectre.  
            \textit{Exemple :} en radio, on translate un signal audio vers une fréquence porteuse.

            \item \textbf{Changement d'échelle temporelle :}
            \[
            x(at) \longleftrightarrow \frac{1}{|a|}X\left(\frac{f}{a}\right)
            \]
            \textit{Signification :} compression temporelle $\Rightarrow$ élargissement du spectre.  
            \textit{Exemple :} un signal très court contient beaucoup de fréquences.

            \item \textbf{Convolution :}
            \[
            x(t)*h(t) \longleftrightarrow X(f)H(f)
            \]
            \textit{Signification :} filtrer un signal revient à multiplier son spectre.  
            \textit{Exemple :} un filtre passe-bas supprime les hautes fréquences.

            \item \textbf{Produit temporel :}
            \[
            x(t)h(t) \longleftrightarrow X(f)*H(f)
            \]
            \textit{Signification :} multiplier deux signaux mélange leurs fréquences.  
            \textit{Exemple :} modulation d’amplitude (AM).

        \end{itemize}

        \textbf{SLI en fréquence :}
        \[
        Y(f) = H(f) X(f)
        \]
    
    \end{multicols}
    
    \begin{multicols}{2}[
    \section{Échantillonnage}
    ]

    \textbf{Principe :} transformation d'un signal continu en signal discret.

    \textbf{Fréquence d'échantillonnage :} $f_e$

    \textbf{Théorème de Nyquist-Shannon :}
    \[
    f_e \ge 2 f_{\max}
    \]

    Sinon : \textbf{repliement de spectre (aliasing)}

    \textbf{Signal échantillonné :}
    \[
    x_e(t) = x(t) \sum_{n=-\infty}^{\infty} \delta(t-nT_e)
    \]

    \section{Filtrage des signaux}

    \textbf{But :} modifier le contenu fréquentiel d'un signal.

    \textbf{Types de filtres :}
    \begin{itemize}
        \item Passe-bas
        \item Passe-haut
        \item Passe-bande
        \item Coupe-bande
    \end{itemize}

    \textbf{Réponse fréquentielle :} $H(f)$

    \textbf{Filtrage :}
    \begin{itemize}
        \item Temporel : convolution
        \item Fréquentiel : multiplication
    \end{itemize}

    \textbf{Exemple : filtre passe-bas idéal}
    \[
    H(f)=
    \begin{cases}
    1 & |f| < f_c \\
    0 & \text{sinon}
    \end{cases}
    \]
    
    \end{multicols}

    \begin{multicols}{2}[
    \section{Résumé}
    ]

    \begin{itemize}
        \item Signal = information temporelle
        \item SLI = convolution
        \item Fourier = analyse fréquentielle
        \item Échantillonnage = discretisation
        \item Filtrage = sélection fréquentielle
    \end{itemize}
    
    \end{multicols}

    \section{Transformées de Fourier usuelles}

    \textbf{Convention :} $X(f) = \int x(t)e^{-j2\pi ft}dt$

    \begin{center}
    \begin{tabular}{|c|c|}
    \hline
    \textbf{Signal temporel} $x(t)$ & \textbf{Transformée} $X(f)$ \\
    \hline

    $\delta(t)$ & $1$ \\
    \hline

    $1$ & $\delta(f)$ \\
    \hline

    $e^{j2\pi f_0 t}$ & $\delta(f - f_0)$ \\
    \hline

    $\cos(2\pi f_0 t)$ & $\frac{1}{2}\left[\delta(f - f_0) + \delta(f + f_0)\right]$ \\
    \hline

    $\sin(2\pi f_0 t)$ & $\frac{1}{2j}\left[\delta(f - f_0) - \delta(f + f_0)\right]$ \\
    \hline

    $\text{rect}(t/T)$ & $T \cdot \text{sinc}(Tf)$ \\
    \hline

    $\text{sinc}(t)$ & $\text{rect}(f)$ \\
    \hline

    $e^{-at}u(t),\ a>0$ & $\frac{1}{a + j2\pi f}$ \\
    \hline

    \end{tabular}
    \end{center}

    \vspace{0.2cm}

    \textbf{Définitions utiles :}
    \[
    \text{sinc}(t) = \frac{\sin(\pi t)}{\pi t}
    \]

    

    \textbf{À retenir :}
    \begin{itemize}
        \item Impulsion $\leftrightarrow$ constante
        \item Sinusoïde $\leftrightarrow$ pics de Dirac
        \item Signal court $\leftrightarrow$ spectre large
        \item Rect $\leftrightarrow$ sinc (dualité temps/fréquence)
    \end{itemize}

\end{document}